% !TeX program = pdflatex
\documentclass[11pt,a4paper]{article}

\usepackage[utf8]{inputenc}
\usepackage[T1]{fontenc}
\usepackage[spanish]{babel}
\usepackage{amsmath,amssymb}
\usepackage{geometry}
\usepackage{graphicx}
\usepackage{booktabs}
\usepackage{siunitx}
\usepackage{hyperref}
\usepackage{float}

\geometry{margin=1.5cm}
\graphicspath{{figures/}}

\title{TEMAS SELECTOS DE C\'OMPUTO DE ALTO RENDIMIENTO PARA SIMULACIONES F\'ISICAS\\
Analisis Num\'erico y Computacional\\
Ecuacion de Calor/Laplace en 1D y 2D\\
Dr. Juan Carlos Cajas}
\author{Autor: Emilio Ramirez}
\date{\today}

\begin{document}

\maketitle

\begin{abstract}
Este documento presenta la formulaci\'on matem\'atica, discretizaci\'on por diferencias
finitas y evaluacion computacional de solucionadores para ecuaciones de tipo
calor/Laplace. Se usa el algoritmo de Thomas para sistemas tridiagonales y se
analiza el rendimiento de la paralelizaci\'on con OpenMP en el caso bidimensional.
\end{abstract}

\section{Introducci\'on}
El objetivo es documentar una linea base reproducible para comparar variantes de
implementacion sin cambiar el paradigma matem\'atico ni de paralelizaci\'on.

\section{Modelo matem\'atico}
\subsection{Caso 1D estacionario}
\begin{equation}
\frac{d}{dx}\left(k\frac{dT}{dx}\right) + \dot{q} = 0
\end{equation}

Para $k$ constante (material isotr\'opico y $\Delta T$ peque\~nos) 
y sin fuentes internas ($\dot{q}=0$), se obtiene la Ecuaci\'on de Laplace:

\begin{equation}
\frac{d^2T}{dx^2}=0
\end{equation}

\subsection{Caso 2D estacionario}

Para el caso bidimensional estacionario, se obtiene la ecuaci\'on de Laplace en 2D:

\begin{equation}
\frac{\partial}{\partial x}\left(k\frac{\partial T}{\partial x}\right) + \frac{\partial}{\partial y}\left(k\frac{\partial T}{\partial y}\right) + \dot{q} = 0
\end{equation}

\begin{equation}
\Rightarrow \frac{\partial^2T}{\partial x^2}+\frac{\partial^2T}{\partial y^2}=0
\end{equation}

\begin{equation}
\Rightarrow \nabla^2 T = 0
\end{equation}

\subsection{Condiciones de frontera}

Se consideran condiciones de frontera de tipo Dirichlet (temperatura fija) y 
Neumann (flujo de calor especificado) seg\'un el experimento.

\textbf{Frontera tipo Dirichlet:}

\begin{equation}
T(0)=T_0, \quad T(L)=T_L
\end{equation}

\textbf{Frontera tipo Neumann:}

Si $q_N$ representa el flujo de calor saliente normal a la frontera, la
condici\'on de Neumann se escribe como:

\begin{equation}
-k\nabla T \cdot \mathbf{n} = q_N.
\end{equation}

En el caso unidimensional, con normal exterior $\mathbf{n}=-1$ en $x=0$ y
$\mathbf{n}=1$ en $x=L$, esto equivale a:

\begin{equation}
k\frac{dT}{dx}\bigg|_{x=0} = q_0, \quad
-k\frac{dT}{dx}\bigg|_{x=L} = q_L.
\end{equation}

\section{Discretizacion del caso 1D}

La ecuaci\'on de Laplace se discretiza usando diferencias finitas.

\subsection{Ecuaci\'on interior}

Se parte de la ecuaci\'on diferencial Eq. (2) discretizada:

\begin{equation}
\frac{\Delta}{\Delta x} \left( \frac{\Delta T}{\Delta x} \right) = 0
\end{equation}

\begin{equation}
\Rightarrow \frac{\Delta}{\Delta x} \left( \frac{T(i + 1) - T(i)}{\Delta x} \right) = 0
\end{equation}

\begin{equation}
\Rightarrow\frac{ \Delta T(i + 1) - \Delta T(i)}{ (\Delta x)^2} = 0
\end{equation}

\begin{equation}
\Rightarrow \frac{\left( T(i + 2) - T(i + 1) \right) - 
\left( T(i + 1) - T(i) \right)}{ (\Delta x)^2} = 0
\end{equation}

\begin{equation}
\Rightarrow T_{i+2} - 2 T_{i+1} + T_{i} = 0
\end{equation}

Cambiando el \'indice $i+1 \to i$:

\begin{equation}
T_{i+1} -2T_i + T_{i-1}=0
\end{equation}

con condiciones de frontera de tipo Dirichlet/Neumann seg\'un el experimento.

\textbf{Ecuaci\'on matricial 1D:}

La Eq. (13) se puede escribir en forma matricial $A \mathbf{T} = \mathbf{r}$,
donde $A$ es una matriz tridiagonal con $-2$ en la diagonal principal y $1$ 
en las diagonales adyacentes, $\mathbf{T}$ es el vector de temperaturas 
desconocidas y $\mathbf{r}$ es el vector de condiciones de frontera.

\begin{equation}
\begin{pmatrix}
1 & -2 & 1 & 0 & 0  \cdots & 0 & 0 & 0 \\
0 & 1 & -2 & 1 & 0  \cdots & 0 & 0 & 0 \\
0 & 0 & 1 & -2 & 1  \cdots & 0 & 0 & 0 \\
\vdots & \vdots & \vdots & \ddots & \vdots \\
0 & 0 & 0 & 0 & 0  \cdots & 1 & -2 & 1
\end{pmatrix}
\begin{pmatrix}
T_2 \\
T_3 \\
T_4 \\
\vdots \\
T_{n-1}
\end{pmatrix}
=
\begin{pmatrix}
0 \\
0 \\
0 \\
\vdots \\
0
\end{pmatrix}
\end{equation}

Se hace notar que faltan las temperaturas frontera i.e. $T_1$ y $T_n$ que 
depender\'an del tipo de frontera (Dirichlet o Neumann).

\textbf{Ecuaci\'on matricial 1D con condiciones frontera tipo Dirichlet:}

Con condiciones de frontera de tipo Dirichlet, el sistema se modifica para 
incluir los valores conocidos en las matrices. Dado que $T_1$ y $T_n$ son 
conocidos:

\begin{equation}
T_1 = T_0, \quad T_n = T_L
\end{equation}

\begin{equation}
\begin{pmatrix}
1 & 0 & 0 & 0 & 0  \cdots & 0 & 0 & 0 \\
1 & -2 & 1 & 0 & 0  \cdots & 0 & 0 & 0 \\
0 & 1 & -2 & 1 & 0  \cdots & 0 & 0 & 0 \\
\vdots & \vdots & \vdots & \ddots & \vdots \\
0 & 0 & 0 & 0 & 0  \cdots & 1 & -2 & 1 \\
0 & 0 & 0 & 0 & 0  \cdots & 0 & 0 & 1
\end{pmatrix}
\begin{pmatrix}
T_1 \\
T_2 \\
T_3 \\
\vdots \\
T_{n-1} \\
T_n
\end{pmatrix}
=
\begin{pmatrix}
T_0 \\
0 \\
0 \\
\vdots \\
0 \\
T_L
\end{pmatrix}
\end{equation}

\textbf{Ecuaci\'on matricial 1D con frontera izquierda tipo Neumann:}

Se estudia el caso en el que la frontera izquierda usa una condici\'on de Neumann
y la frontera derecha conserva una condici\'on de Dirichlet. Con flujo saliente
normal $q_0$ en $x=0$:

\begin{equation}
k\frac{dT}{dx}\bigg|_{x=0}=q_0.
\end{equation}

Aproximando la derivada normal con una diferencia hacia adelante:

\begin{equation}
\frac{dT}{dx}\bigg|_{x=0} \approx \frac{T_2-T_1}{\Delta x},
\end{equation}

se obtiene:

\begin{equation}
k\frac{T_2-T_1}{\Delta x}=q_0
\quad \Rightarrow \quad
-T_1+T_2=\frac{q_0}{k}\Delta x.
\end{equation}

Por tanto, el sistema algebraico para una frontera izquierda Neumann y una
frontera derecha Dirichlet $T_n=T_L$ queda:

\begin{equation}
\begin{pmatrix}
-1 & 1 & 0 & 0 & 0  \cdots & 0 & 0 & 0 \\
1 & -2 & 1 & 0 & 0  \cdots & 0 & 0 & 0 \\
0 & 1 & -2 & 1 & 0  \cdots & 0 & 0 & 0 \\
\vdots & \vdots & \vdots & \ddots & \vdots \\
0 & 0 & 0 & 0 & 0  \cdots & 1 & -2 & 1 \\
0 & 0 & 0 & 0 & 0  \cdots & 0 & 0 & 1
\end{pmatrix}
\begin{pmatrix}
T_1 \\
T_2 \\
T_3 \\
\vdots \\
T_{n-1} \\
T_n
\end{pmatrix}
=
\begin{pmatrix}
\dfrac{q_0}{k}\Delta x \\
0 \\
0 \\
\vdots \\
0 \\
T_L
\end{pmatrix}.
\end{equation}

Si se imponen condiciones de Neumann en ambos extremos para la ecuaci\'on de
Laplace sin fuentes, el sistema queda determinado solo hasta una constante
aditiva y requiere una condici\'on adicional de referencia. Adem\'as, los flujos
deben cumplir la condici\'on de compatibilidad $q_0+q_L=0$.

\section{Algoritmo de Thomas}

Tenemos el sistema tridiagonal $A \mathbf{T} = \mathbf{r}$, donde $A$ es 
una matriz tridiagonal con elementos $a_i$ en la subdiagonal, $b_i$ en la diagonal,
y $c_i$ en la superdiagonal. De la misma manera, $\mathbf{T}$ es el vector de 
temperaturas desconocidas y $\mathbf{r}$ es el vector de condiciones de frontera,
con elementos $d_i$.

\begin{align*}
b_1 T_1 + c_1 T_2 &= r_1 \\
a_2 T_1 + b_2 T_2 + c_2 T_3 &= r_2 \\
a_3 T_2 + b_3 T_3 + c_3 T_4 &= r_3 \\
&\vdots \\
a_{n-1} T_{n-2} + b_{n-1} T_{n-1} + c_{n-1} T_n &= r_{n-1} \\
a_n T_{n-1} + b_n T_n &= r_n
\end{align*}

\begin{equation}
\begin{pmatrix}
b_1 & c_1 & 0 & 0 & 0  \cdots & 0 & 0 & 0 \\
a_2 & b_2 & c_2 & 0 & 0  \cdots & 0 & 0 & 0 \\
0 & a_3 & b_3 & c_3 & 0  \cdots & 0 & 0 & 0 \\
\vdots & \vdots & \vdots & \ddots & \vdots \\
0 & 0 & 0 & 0 & 0  \cdots & a_{n-1} & b_{n-1} & c_{n-1} \\
0 & 0 & 0 & 0 & 0  \cdots & 0 & a_n & b_n
\end{pmatrix}
\begin{pmatrix}T_1 \\
T_2 \\
T_3 \\
\vdots \\
T_{n-1} \\
T_n
\end{pmatrix}
=
\begin{pmatrix}r_1 \\
r_2 \\
r_3 \\
\vdots \\
r_{n-1} \\
r_n
\end{pmatrix}
\end{equation}

As\'i, se tienen las ecuaciones del sistema tridiagonal a resolver:

\begin{equation}
a_i T_{i-1}+b_iT_i+c_iT_{i+1}=r_i,
\quad i=1,2,\ldots,n, \quad a_1=0, \quad c_n=0.
\end{equation}

Resolviendo,

\begin{equation}
T_1 = \frac{r_1 - c_1 T_2}{b_1}
\end{equation}

Y definimos,

\begin{equation}
\tilde{r_1} = \frac{r_1}{b_1}, \quad \tilde{c_1} = \frac{c_1}{b_1}
\end{equation}

Con lo cual, la primera ecuaci\'on se reescribe como:

\begin{equation}
T_1 = \tilde{r_1} - \tilde{c_1} T_2
\end{equation}

Si ahora resolvemos el segundo rengl\'on para $T_2$:

\begin{equation}
T_2 = \frac{r_2 - a_2 T_1 - c_2 T_3}{b_2}
\end{equation}

\begin{equation}
\Rightarrow T_2 = \frac{r_2 - a_2 (\tilde{r_1} - \tilde{c_1} T_2) - c_2 T_3}{b_2}
\end{equation}

\begin{equation}
\Rightarrow T_2 \left(1 - \frac{a_2 \tilde{c_1}}{b_2}\right) = \frac{r_2 - a_2 \tilde{r_1} - c_2 T_3}{b_2}
\end{equation}

\begin{equation}
\Rightarrow T_2 = \frac{r_2 - a_2 \tilde{r_1} - c_2 T_3}{b_2 - a_2 \tilde{c_1}}
\end{equation}

Y si ahora definimos 

\begin{equation}
\tilde{r_2} = \frac{r_2 - a_2 \tilde{r_1}}{b_2 - a_2 \tilde{c_1}}, \quad 
\tilde{c_2} = \frac{c_2}{b_2 - a_2 \tilde{c_1}}
\end{equation}

Entonces nos queda,

\begin{equation}
T_2 = \tilde{r_2} - \tilde{c_2} T_3
\end{equation}

En general, para $i=2,3,\ldots,n-1$ se tiene:

\begin{equation}
T_i = \tilde{r_i} - \tilde{c_i} T_{i+1},
\end{equation}

con:

\begin{equation}
\tilde{r_i} = \frac{r_i - a_i \tilde{r}_{i-1}}{b_i - a_i \tilde{c}_{i-1}}, \quad
\tilde{c_i} = \frac{c_i}{b_i - a_i \tilde{c}_{i-1}}.
\end{equation}

Para $T_n$ se tiene:

\begin{equation}
T_n = \frac{r_n - a_n T_{n-1}}{b_n}
\end{equation}

\begin{equation}
\Rightarrow T_n = \frac{r_n - a_n 
\left( \tilde{r}_{n-1} - \tilde{c}_{n-1} T_{n} \right)}{b_n}
\end{equation}

\begin{equation}
\Rightarrow T_n \left(1 + \frac{a_n \tilde{c}_{n-1}}{b_n}\right) 
= \frac{r_n - a_n \tilde{r}_{n-1}}{b_n}
\end{equation}

\begin{equation}
\Rightarrow T_n = \frac{r_n - a_n \tilde{r}_{n-1}}{b_n + a_n \tilde{c}_{n-1}}
\end{equation}

\begin{equation}
\Rightarrow T_n = \tilde{r}_n
\end{equation}

\section{Metodologia computacional del caso 1D}

Para el caso unidimensional se resuelve directamente el sistema tridiagonal
producido por la discretizaci\'on de Laplace:

\begin{equation}
A\mathbf{T}=\mathbf{r}.
\end{equation}

Se implementa el algoritmo de Thomas para resolver este sistema de forma 
eficiente. Se usa el comando \verb|gfortran tridiagonal.f90 calor1D.f90 -o calor1D| 
para compilar la versi\'on secuencial. El programa se ejecuta con \verb|./calor1D| 
y los resultados se guardan en archivos de texto para su posterior an\'alisis.

\medskip{}

Se muestra a continuaci\'on el c\'odigo Fortran de la versi\'on secuencial 
para el caso 1D. Se comienza abriendo el programa y declarando las variables 
necesarias para el algoritmo:

\begin{verbatim}
Program Calor1D
  Implicit none
  integer :: ii
  integer, parameter :: nn = 60
  integer, parameter :: L = 10 ! metros
  double precision :: deltax
  double precision, parameter :: q0_sobre_k = -1.d0
  double precision :: tt(nn)    ! vector de incognitas
  double precision :: rr(nn)    ! vector de resultados del s. ecuaciones
  double precision :: aa(nn), bb(nn), cc(nn) ! Variables para almacenar
  double precision :: cp(nn), den(nn)
\end{verbatim}

Establecemos el dominio computacional y el paso de malla $\Delta x = L/nn$:
El valor $L=10$ metros define la longitud total del dominio de soluci\'on; se elige
arbitrariamente para este problema de prueba. La discretizaci\'on se realiza 
con $nn=60$ nodos, resultando en $\Delta x = L/60$.

\begin{verbatim}
  ! Dominio computacional
  deltax  = dble(L)/nn
\end{verbatim}

Luego, se inicializan los vectores necesarios para el algoritmo de Thomas.
Se hace notar que no se construye la matriz completa $A$, sino que se almacenan 
solamente las diagonales $a$, $b$ y $c$ en los vectores correspondientes.

\begin{verbatim}
  ! Inicializacion de variables
  !
  aa(:)   = 0.d0
  bb(:)   = 0.d0
  cc(:)   = 0.d0
  rr(:)   = 0.d0
\end{verbatim}

Se ensamblan los coeficientes de la matriz tridiagonal y el vector de resultados
seg\'un la discretizaci\'on de la ecuaci\'on de Laplace. Para los nodos 
interiores, los coeficientes son constantes: $a_i=1$, $b_i=-2$ y $c_i=1$.

\begin{verbatim}
  ! Ensamblar la matriz tridiagonal y el vector de resultados
  do ii = 2, nn-1
     !
     aa(ii) = 1.d0
     bb(ii) =-2.d0
     cc(ii) = 1.d0
     rr(ii) = 0.d0
     !
  end do
\end{verbatim}

A continuaci\'on, se imponen las condiciones de frontera. En este caso, 
se asume una condici\'on de Neumann en el extremo izquierdo ($x=0$) con 
flujo $q_0$ y una condici\'on de Dirichlet en el extremo derecho ($x=L$) 
con temperatura $T_L$. Esto se traduce en los coeficientes de la primera y 
\'ultima ecuaciones del sistema.

\begin{verbatim}
  ! Impone cond. frontera
  !
  aa(1)     = 0.d0 ! no se usa en los c'alculos
  bb(1)     =-1.d0
  cc(1)     = 1.d0
  rr(1)     = q0_sobre_k * deltax
  !  
  aa(nn)    = 0.d0
  bb(nn)    = 1.d0
  cc(nn)    = 0.d0 ! no se usa en los c'alculos
  rr(nn)    = 0.d0
\end{verbatim}

Se resuelve el sistema tridiagonal usando el algoritmo de Thomas, que se 
implementa en las subrutinas \verb|tri_factor| y \verb|tri_solve|. Estas
subrutinas realizan la factorizaci\'on y la resoluci\'on del sistema de 
forma eficiente, aprovechando la estructura tridiagonal de la matriz.

\begin{verbatim}
  ! Resolver el problema algebraico
  !
  call tri_factor(aa,bb,cc,cp,den,nn)
  call tri_solve(aa,cp,den,rr,tt,nn)
\end{verbatim}

Finalmente, se guarda el resultado en un archivo de texto para su posterior
an\'alisis. Se escribe el vector de temperaturas $tt$ junto con las posiciones
correspondientes en el dominio.

\begin{verbatim}
  ! Postproceso
  !
  do ii = 1, nn
     write(101,*) (ii-1)*deltax, tt(ii)
  end do
  !
end Program Calor1D
\end{verbatim}

\textbf{Subrutinas: \texttt{tri, tri\_solver y tri\_factor}:}  

La subrutina \verb|tri| resuelve el sistema tridiagonal completo en un solo
paso: primero realiza la eliminaci\'on gaussiana hacia adelante y luego la
sustituci\'on hacia atr\'as.

\medskip

\textit{Barrido hacia adelante (\texttt{gausselim}):}
Para $i = 2, 3, \ldots, n$, se calcula el multiplicador de eliminaci\'on y
se actualizan la diagonal y el lado derecho \textit{in-place}:

\begin{align}
m_i          &= \frac{a_i}{b_{i-1}}, \\
b_i          &\leftarrow b_i - m_i\, c_{i-1}, \\
r_i          &\leftarrow r_i - m_i\, r_{i-1}.
\end{align}

Tras este barrido el sistema original $A\mathbf{u}=\mathbf{r}$ ha quedado
reducido a uno triangular superior equivalente $U\mathbf{u}=\mathbf{r}^*$,
donde $\mathbf{r}^*$ denota el lado derecho transformado.

\medskip

\textit{Sustituci\'on hacia atr\'as (\texttt{sustatras}):}
Se obtiene la \'ultima inc\'ognita directamente y luego se retrocede:

\begin{align}
u_n &= \frac{r_n^{\,*}}{b_n^{\,*}}, \\[4pt]
u_i &= \frac{r_i^{\,*} - c_i\, u_{i+1}}{b_i^{\,*}},
\quad i = n-1,\, n-2,\, \ldots,\, 1.
\end{align}

El coste total es $\mathcal{O}(n)$ operaciones, frente a los
$\mathcal{O}(n^3)$ de una eliminaci\'on gaussiana general.

\textbf{Subrutina tri\_factor:}

Se explica a continuaci\'on el c\'odigo de la subrutina \verb|tri_factor|, 
que realiza la factorizaci\'on de Thomas para una matriz tridiagonal 
constante.

\medskip{}

Se declaran las variables de entrada y salida, incluyendo los vectores de 
coeficientes de la matriz tridiagonal y los vectores para almacenar la 
factorizaci\'on.

\begin{verbatim}
  subroutine tri(a,b,c,r,u,n)
  !
  ! Subrutina que resuelve el problema
  !
  ! Ax=r donde A es una matriz tridiagonal con
  ! elementos a, b, c
  ! y devuelve el valor de u
  !
  implicit none
  !
  integer,          intent(in)    :: n
  double precision, intent(in)    :: a(n),c(n)
  double precision, intent(inout) :: b(n),r(n)
  double precision, intent(inout) :: u(n)
  integer :: i
  double precision :: m

  ! eliminacion elementos bajo la matriz
  gausselim: do i=2,n
     m = a(i)/b(i-1)
     r(i)=r(i)-m*r(i-1)
     b(i)=b(i)-m*c(i-1)
  end do gausselim
  !
  ! solucion para u
  !
  u(n)=r(n)/b(n)
  sustatras: do i=n-1,1,-1
     u(i)=(r(i)-c(i)*u(i+1))/b(i)
  end do sustatras
  !
end subroutine tri

subroutine tri_factor(a,b,c,cp,den,n)
  !
  ! Factorizacion de Thomas para una matriz tridiagonal constante.
  !
  implicit none
  !
  integer,          intent(in)  :: n
  double precision, intent(in)  :: a(n),b(n),c(n)
  double precision, intent(out) :: cp(n),den(n)
  integer :: i

  den(1) = b(1)
  cp(1)  = c(1)/den(1)

  do i=2,n-1
     den(i) = b(i) - a(i)*cp(i-1)
     cp(i) = c(i)/den(i)
  end do
  den(n) = b(n) - a(n)*cp(n-1)
  cp(n) = 0.d0
end subroutine tri_factor

subroutine tri_solve(a,cp,den,r,u,n)
  !
  ! Sustitucion usando la factorizacion calculada por tri_factor.
  ! Modifica r in-place para guardar el lado derecho transformado.
  !
  implicit none
  !
  integer,          intent(in)    :: n
  double precision, intent(in)    :: a(n),cp(n),den(n)
  double precision, intent(inout) :: r(n)
  double precision, intent(inout) :: u(n)
  integer :: i

  r(1) = r(1)/den(1)
  do i=2,n
     r(i) = (r(i) - a(i)*r(i-1))/den(i)
  end do

  u(n) = r(n)
  do i=n-1,1,-1
     u(i) = r(i) - cp(i)*u(i+1)
  end do
end subroutine tri_solve

\end{verbatim}


\section{Discretizaci\'on del caso 2D}

En el caso bidimensional, la ecuaci\'on de Laplace se discretiza como:

\begin{equation}
\frac{T_{i+1,j}-2T_{i,j}+T_{i-1,j}}{(\Delta x)^2}
+
\frac{T_{i,j+1}-2T_{i,j}+T_{i,j-1}}{(\Delta y)^2}
=0.
\end{equation}

Separando los t\'erminos en la direcci\'on $x$:

\begin{equation}
\frac{1}{(\Delta x)^2}T_{i-1,j}
-2\left(\frac{1}{(\Delta x)^2}+\frac{1}{(\Delta y)^2}\right)T_{i,j}
+\frac{1}{(\Delta x)^2}T_{i+1,j}
=-\frac{T_{i,j-1}+T_{i,j+1}}{(\Delta y)^2}.
\end{equation}

Esta expresi\'on genera un sistema tridiagonal por cada linea horizontal
$j$:

\begin{equation}
a_i^{(x)}T_{i-1,j}+b_i^{(x)}T_{i,j}+c_i^{(x)}T_{i+1,j}=r_i^{(x)}.
\end{equation}

Los coeficientes usados son:

\begin{align}
a_i^{(x)} &= \frac{1}{(\Delta x)^2}, \\
b_i^{(x)} &= -2\left(\frac{1}{(\Delta x)^2}+\frac{1}{(\Delta y)^2}\right), \\
c_i^{(x)} &= \frac{1}{(\Delta x)^2}, \\
r_i^{(x)} &= -\frac{T_{i,j-1}+T_{i,j+1}}{(\Delta y)^2}.
\end{align}

De forma an\'aloga, separando los t\'erminos en la direcci\'on $y$:

\begin{equation}
\frac{1}{(\Delta y)^2}T_{i,j-1}
-2\left(\frac{1}{(\Delta x)^2}+\frac{1}{(\Delta y)^2}\right)T_{i,j}
+\frac{1}{(\Delta y)^2}T_{i,j+1}
=-\frac{T_{i-1,j}+T_{i+1,j}}{(\Delta x)^2}.
\end{equation}

Esto produce sistemas tridiagonales verticales:

\begin{equation}
a_j^{(y)}T_{i,j-1}+b_j^{(y)}T_{i,j}+c_j^{(y)}T_{i,j+1}=r_j^{(y)},
\end{equation}

con:

\begin{align}
a_j^{(y)} &= \frac{1}{(\Delta y)^2}, \\
b_j^{(y)} &= -2\left(\frac{1}{(\Delta x)^2}+\frac{1}{(\Delta y)^2}\right), \\
c_j^{(y)} &= \frac{1}{(\Delta y)^2}, \\
r_j^{(y)} &= -\frac{T_{i-1,j}+T_{i+1,j}}{(\Delta x)^2}.
\end{align}

El proceso iterativo se detiene cuando el cambio entre dos iteraciones
consecutivas es menor que una tolerancia prescrita. El residuo usado es:

\begin{equation}
R^{(k)}=
\left[
\sum_{i=1}^{n_x}\sum_{j=1}^{n_y}
\left(T_{i,j}^{(k)}-T_{i,j}^{(k-1)}\right)^2
\right]^{1/2}.
\end{equation}

El criterio de convergencia es:

\begin{equation}
R^{(k)}<\varepsilon.
\end{equation}

\section{Metodologia computacional secuencial del caso 2D}

El algoritmo secuencial para el caso 2D se implementa en el programa 
\texttt{secuencial\_calor1D\_calor2D.f90}. Se resuelve iterativamente el 
sistema de ecuaciones generado por la discretización bidimensional de la 
ecuación de Laplace. Se usa el mismo algoritmo de Thomas para resolver los 
sistemas tridiagonales en cada dirección. El programa se ejecuta con 
\verb|./calor2D| y los resultados se guardan en archivos de texto para su 
posterior análisis.

La versi\'on paralela distribuye los barridos independientes entre hilos de
OpenMP. Si $t_p$ es el tiempo de ejecuci\'on usando $p$ hilos y $t_1$ es el
tiempo de referencia secuencial, el speedup se calcula como:

\begin{equation}
S_p=\frac{t_1}{t_p}.
\end{equation}

La eficiencia paralela mide qu\'e fracci\'on del uso ideal de los $p$ hilos se
aprovecha:

\begin{equation}
E_p=\frac{S_p}{p}.
\end{equation}

El caso ideal corresponde a:

\begin{equation}
S_p^{ideal}=p,
\quad
E_p^{ideal}=1.
\end{equation}

\section{Resultados}

Las tablas de rendimiento se guardan en \texttt{docs/tables/} y las figuras en
\texttt{docs/figures/}. Para cada n\'umero de hilos $p$, el tiempo reportado se
puede obtener como el promedio de varias corridas:

\begin{equation}
\overline{t}_p=\frac{1}{N}\sum_{\ell=1}^{N}t_{p,\ell},
\end{equation}

donde $N$ es el n\'umero de repeticiones y $t_{p,\ell}$ es el tiempo medido en
la corrida $\ell$ usando $p$ hilos. Con este promedio, el speedup experimental
queda:

\begin{equation}
S_p^{real}=\frac{\overline{t}_1}{\overline{t}_p}.
\end{equation}

La eficiencia experimental se calcula como:

\begin{equation}
E_p^{real}=\frac{S_p^{real}}{p}
=\frac{\overline{t}_1}{p\,\overline{t}_p}.
\end{equation}

Para comparar con la aceleraci\'on ideal se usa:

\begin{equation}
\Delta S_p = S_p^{ideal}-S_p^{real}=p-S_p^{real}.
\end{equation}

Una diferencia $\Delta S_p$ cercana a cero indica que la ejecuci\'on paralela se
aproxima al comportamiento ideal. Cuando $\Delta S_p$ crece, la eficiencia baja
por costos de sincronizaci\'on, acceso a memoria, regiones no paralelas o
desbalance entre hilos.

\subsection{Comparaci\'on final con el c\'odigo de referencia}

Se compar\'o la versi\'on paralela desarrollada en \texttt{paralelo/src/} contra
el c\'odigo de referencia proporcionado en \texttt{paralelo/reference/}. Los
archivos de referencia no fueron modificados. Para reducir el efecto de
arranques fr\'ios y procesos de fondo, se us\'o una corrida de calentamiento por
n\'umero de hilos y se report\'o la mediana de seis corridas medidas.

Para aislar el efecto del paralelismo bajo el mismo comportamiento num\'erico,
en la versi\'on desarrollada se reprodujeron dos errores presentes en
\texttt{reference/}: (i) la condici\'on de frontera inferior del barrido en
direcci\'on $y$ se dej\'o con la misma formulaci\'on incorrecta, y (ii) el corte
por convergencia basado en residuo se mantuvo desactivado, forzando la
ejecuci\'on hasta \texttt{itermax}.

La comparaci\'on debe leerse como una comparaci\'on de comportamiento paralelo en
la misma m\'aquina, no como igualdad estricta de carga num\'erica: ambas versiones 
usan una malla de $480\times 240$.

\begin{table}[H]
\centering
\caption{Comparaci\'on de tiempos y aceleraci\'on en la misma m\'aquina.}
\label{tab:comparacion-speedup-final}
\begin{tabular}{rrrrr}
\toprule
Hilos & $t_{\mathrm{nuestra}}$ [s] & $S_{\mathrm{nuestra}}$ &
$t_{\mathrm{ref}}$ [s] & $S_{\mathrm{ref}}$ \\
\midrule
1 & 40.466 & 1.000 & 36.954 & 1.000 \\
2 & 21.829 & 1.854 & 21.855 & 1.691 \\
3 & 15.465 & 2.617 & 15.736 & 2.348 \\
4 & 13.534 & 2.990 & 14.136 & 2.614 \\
5 & 12.845 & 3.150 & 10.911 & 3.387 \\
6 & 11.342 & 3.568 &  9.832 & 3.759 \\
7 &  8.554 & 4.731 &  8.890 & 4.157 \\
8 &  8.102 & 4.994 & 10.386 & 3.558 \\
\bottomrule
\end{tabular}
\end{table}

\begin{figure}[H]
\centering
\includegraphics[width=0.78\textwidth]{comparacion_speedup_final.png}
\caption{Curvas de aceleraci\'on de la versi\'on desarrollada y la referencia.}
\label{fig:comparacion-speedup-final}
\end{figure}


\section{Conclusiones}



\bibliographystyle{plain}
\bibliography{bib/references}

\end{document}
